\documentclass[11pt,a4paper]{article}

\usepackage[utf8]{inputenc}
\usepackage[T1]{fontenc}
\usepackage[spanish]{babel}
\usepackage{amsmath,amssymb,mathtools}
\usepackage{geometry}
\usepackage{lmodern}
\usepackage{microtype}

\geometry{margin=2.4cm}
\setlength{\parindent}{0pt}
\setlength{\parskip}{0.45em}

\newcommand{\R}{\mathbb{R}}
\newcommand{\sen}{\operatorname{sen}}

\title{Soluciones Tema 1: Conceptos básicos de cálculo vectorial}
\author{Juan Rodriguez}
\date{}

\begin{document}

\maketitle

\section{Ejercicio 1}

Obtén la forma implícita de la curva cuya parametrización es

\[
\alpha(t)=(t^2-4,4t).
\]

\subsection{Solución}

La parametrización es:

\[
\alpha(t)=(t^2-4,4t).
\]

\[
x(t) = t^2-4 \quad y(t) = 4t 
\]

De:

\[
y=4t
\]

se obtiene:

\[
t=\frac{y}{4}.
\]

\[
x=\frac{y^2}{16}-4.
\]

Por tanto:

\[
\boxed{y^2=16(x+4)}
\]

\section{Ejercicio 2}

Determina la expresión implícita de la curva parametrizada por

\[
\alpha(t)=(\sec(t),\tan(t)).
\]

\subsection{Solución}

La parametrización es:

\[
\alpha(t)=(\sec t,\tan t).
\]

\[
\text{Usamos: } \quad \boxed{\sec^2t-\tan^2t=1} \rightarrow x^2-y^2=1
\]

\[
\boxed{x^2-y^2=1}
\]

\section{Ejercicio 3}

Determina los puntos dobles de la curva

\[
x^2(x-y)+y^2=0
\]

y analiza de qué tipo son.

\subsection{Solución}

Sea:

\[
F(x,y)=x^2(x-y)+y^2.
\]

Los puntos singulares verifican:

\[
F=F_x=F_y=0.
\]

Se tiene:

\[
F_x=3x^2-2xy,
\]

\[
F_y=-x^2+2y.
\]

La única solución es:

\[
\boxed{(0,0)}
\]

La matriz hessiana en el origen es:

\[
H_F(0,0)
=
\begin{pmatrix}
0&0\\
0&2
\end{pmatrix}
\]

y:

\[
|H_F(0,0)|=0.
\]

Además, la ecuación puede verse como una cuadrática en \(y\):

\[
y^2-x^2y+x^3=0.
\]

Su discriminante es:

\[
\Delta
=
x^4-4x^3
=
x^3(x-4).
\]

Para \(x<0\) suficientemente próximo a \(0\) existen dos ramas reales que llegan al origen con la misma tangente:

\[
y=0.
\]

Por tanto:

\[
\boxed{
(0,0)\text{ es una cúspide}
}
\]

\section{Ejercicio 4}
\subsection{Solución}

Sea:

\[
F(x,y)
=
x^3+x^2+y^2-x-4y+3.
\]

Las derivadas parciales son:

\[
F_x=3x^2+2x-1,
\]

\[
F_y=2y-4.
\]

Resolviendo:

\[
F=F_x=F_y=0
\]

se obtiene:

\[
\boxed{
(-1,2)
}
\]

La matriz hessiana es:

\[
H_F=
\begin{pmatrix}
6x+2&0\\
0&2
\end{pmatrix}.
\]

En el punto singular:

\[
|H_F(-1,2)|
=
(-4)(2)
=
-8<0.
\]

Por tanto:

\[
\boxed{
(-1,2)\text{ es un nodo}
}
\]

\section{Ejercicio 5}
\subsection{Solución}

Sea:

\[
F(x,y)
=
x^3+y^3-3axy,
\qquad
a\neq0.
\]

Se tiene:

\[
F_x=3x^2-3ay,
\]

\[
F_y=3y^2-3ax.
\]

La única solución de:

\[
F=F_x=F_y=0
\]

es:

\[
\boxed{(0,0)}
\]

La matriz hessiana en el origen es:

\[
H_F(0,0)
=
\begin{pmatrix}
0&-3a\\
-3a&0
\end{pmatrix}.
\]

Por tanto:

\[
|H_F(0,0)|
=
-9a^2<0.
\]

Así:

\[
\boxed{
(0,0)\text{ es un nodo}
}
\]

\section{Ejercicio 6}
\subsection{Solución}

Sea:

\[
F(x,y)
=
y^3-x^3-ax^2,
\qquad
a\neq0.
\]

Las derivadas parciales son:

\[
F_x=-3x^2-2ax,
\]

\[
F_y=3y^2.
\]

El único punto singular perteneciente a la curva es:

\[
\boxed{(0,0)}
\]

La parte de menor grado de la ecuación es:

\[
-ax^2=0.
\]

Por tanto, la tangente en el origen es doble:

\[
x=0.
\]

Además, la curva puede escribirse como:

\[
y^3=x^2(x+a),
\]

por lo que existen puntos reales de la curva arbitrariamente próximos al origen.

Así:

\[
\boxed{
(0,0)\text{ es una cúspide}
}
\]

\section{Ejercicio 7}
\subsection{Solución}

Sea:

\[
F(x,y)
=
y^2-2x^2y-x^4y+2x^4.
\]

Las derivadas parciales son:

\[
F_x
=
-4xy-4x^3y+8x^3,
\]

\[
F_y
=
2y-2x^2-x^4.
\]

La única solución de:

\[
F=F_x=F_y=0
\]

es:

\[
\boxed{(0,0)}
\]

La ecuación es cuadrática en \(y\):

\[
y^2-(2x^2+x^4)y+2x^4=0.
\]

Su discriminante es:

\[
\Delta
=
(2x^2+x^4)^2-8x^4
\]

\[
=
x^4(x^4+4x^2-4).
\]

Para \(x\neq0\) suficientemente próximo a \(0\):

\[
x^4+4x^2-4<0.
\]

Por tanto, no existen otros puntos reales de la curva suficientemente próximos al origen.

Así:

\[
\boxed{
(0,0)\text{ es un punto aislado}
}
\]

\section{Ejercicio 8}
\subsection{Solución}

\textbf{a)}

\[
\alpha(t)=(t^3,t^6).
\]

Se tiene:

\[
\alpha'(t)
=
(3t^2,6t^5).
\]

En:

\[
t=0
\]

se anula:

\[
\alpha'(0)=(0,0).
\]

Por tanto:

\[
\boxed{
\alpha\text{ no es regular}
}
\]

Eliminando el parámetro:

\[
y=x^2.
\]

La curva es una parábola y admite, por ejemplo, la parametrización regular:

\[
\boxed{
\widetilde\alpha(s)=(s,s^2)
}
\]

\textbf{b)}

\[
\beta(t)=(t^2,t^3).
\]

Entonces:

\[
\beta'(t)
=
(2t,3t^2).
\]

Así:

\[
\beta'(0)=(0,0),
\]

por lo que:

\[
\boxed{
\beta\text{ no es regular}
}
\]

Eliminando el parámetro:

\[
y^2=x^3.
\]

La curva presenta una cúspide en el origen, por lo que no existe una parametrización regular de toda la curva.

\[
\boxed{
\text{La curva no admite una parametrización regular global}
}
\]

\textbf{c)}

\[
\gamma(t)
=
(\cos^3t,\sen^3t).
\]

Se tiene:

\[
\gamma'(t)
=
(-3\cos^2t\sen t,\,
3\sen^2t\cos t).
\]

La derivada se anula para:

\[
t=0,\frac{\pi}{2},\pi,\frac{3\pi}{2},2\pi.
\]

Por tanto:

\[
\boxed{
\gamma\text{ no es regular}
}
\]

La curva verifica:

\[
\boxed{
x^{2/3}+y^{2/3}=1
}
\]

y es una astroide, con cúspides en:

\[
(\pm1,0),
\qquad
(0,\pm1).
\]

Por tanto:

\[
\boxed{
\text{No admite una parametrización regular global}
}
\]

\section{Ejercicio 9}
\subsection{Solución}

Sea:

\[
\alpha(t)
=
(t^3-3t,t^2-4).
\]

Se tiene:

\[
x'(t)=3t^2-3,
\]

\[
y'(t)=2t.
\]

La tangente es horizontal cuando:

\[
y'(t)=0
\]

y:

\[
x'(t)\neq0.
\]

Así:

\[
t=0
\]

y el punto es:

\[
\boxed{
(0,-4)
}
\]

La tangente es vertical cuando:

\[
x'(t)=0
\]

y:

\[
y'(t)\neq0.
\]

Por tanto:

\[
t=\pm1.
\]

Los puntos son:

\[
\boxed{
(-2,-3),
\qquad
(2,-3)
}
\]

\section{Ejercicio 10}
\subsection{Solución}

La circunferencia unidad puede parametrizarse como:

\[
\alpha(t)
=
(\cos t,\sen t).
\]

Entonces:

\[
\alpha'(t)
=
(-\sen t,\cos t).
\]

Sea:

\[
P=(x_0,y_0)
=
(\cos t_0,\sen t_0).
\]

La recta tangente tiene vector director:

\[
(-y_0,x_0).
\]

Por tanto:

\[
\boxed{
x_0x+y_0y=1
}
\]

\section{Ejercicio 11}
\subsection{Solución}

La curva:

\[
y=x^3
\]

puede parametrizarse como:

\[
\alpha(t)
=
(t,t^3).
\]

El punto:

\[
(1,1)
\]

corresponde a:

\[
t=1.
\]

Se tiene:

\[
\alpha'(t)
=
(1,3t^2).
\]

Por tanto:

\[
\alpha'(1)
=
(1,3).
\]

La recta tangente es:

\[
\boxed{
y-1=3(x-1)
}
\]

La pendiente de la normal es:

\[
-\frac13.
\]

Así:

\[
\boxed{
y-1
=
-\frac13(x-1)
}
\]

\section{Ejercicio 12}
\subsection{Solución}

Una parametrización de la circunferencia es:

\[
\alpha(t)
=
(-1+3\cos t,\,
1+3\sen t),
\qquad
t\in[0,2\pi].
\]

Su derivada es:

\[
\alpha'(t)
=
(-3\sen t,3\cos t).
\]

Por tanto:

\[
\|\alpha'(t)\|
=
3.
\]

La longitud es:

\[
L
=
\int_0^{2\pi}
3\,dt.
\]

Así:

\[
\boxed{
L=6\pi
}
\]

\section{Ejercicio 13}
\subsection{Solución}

La curva es:

\[
y=\sqrt{x^3}=x^{3/2},
\qquad
x\in[0,3].
\]

Tomamos:

\[
\alpha(t)
=
(t,t^{3/2}),
\qquad
t\in[0,3].
\]

Entonces:

\[
\alpha'(t)
=
\left(
1,\frac32\sqrt t
\right).
\]

Por tanto:

\[
L
=
\int_0^3
\sqrt{
1+\frac94t
}\,dt.
\]

Integrando:

\[
L
=
\frac8{27}
\left[
\left(
1+\frac94t
\right)^{3/2}
\right]_0^3.
\]

Así:

\[
\boxed{
L
=
\frac{31\sqrt{31}-8}{27}
}
\]

\section{Ejercicio 14}
\subsection{Solución}

La esfera es:

\[
x^2+y^2+z^2=4
\]

y el plano:

\[
y+z=2.
\]

El centro de la circunferencia de intersección es:

\[
C=(0,1,1).
\]

Tomamos los vectores ortonormales del plano:

\[
u=(1,0,0),
\]

\[
v=
\frac1{\sqrt2}(0,1,-1).
\]

El radio de la circunferencia es:

\[
\sqrt2.
\]

Por tanto, una parametrización es:

\[
\boxed{
\alpha(t)
=
(\sqrt2\cos t,\,
1+\sen t,\,
1-\sen t),
\qquad
t\in[0,2\pi].
}
\]

\section{Ejercicio 15}
\subsection{Solución}

El cilindro es:

\[
(x-1)^2+y^2=1.
\]

Tomamos:

\[
x=1+\cos t,
\qquad
y=\sen t.
\]

La esfera verifica:

\[
x^2+y^2+z^2=4.
\]

Como:

\[
x^2+y^2=2x,
\]

se obtiene:

\[
z^2=4-2x
=
2-2\cos t
=
4\sen^2\frac t2.
\]

Al tratarse de la semiesfera superior:

\[
z\geq0.
\]

Para:

\[
t\in[0,2\pi],
\]

se tiene:

\[
\sen\frac t2\geq0.
\]

Por tanto:

\[
\boxed{
\alpha(t)
=
\left(
1+\cos t,\,
\sen t,\,
2\sen\frac t2
\right),
\qquad
t\in[0,2\pi].
}
\]

\section{Ejercicio 16}
\subsection{Solución}

La parametrización es:

\[
\alpha(t)
=
\left(
t^2-9,\,
\frac{t^3}{3},\,
2t
\right).
\]

El punto:

\[
P=(-9,0,0)
\]

corresponde a:

\[
t=0.
\]

El punto:

\[
Q=(0,9,6)
\]

corresponde a:

\[
t=3.
\]

Se tiene:

\[
\alpha'(t)
=
(2t,t^2,2).
\]

Por tanto:

\[
\|\alpha'(t)\|
=
\sqrt{
4t^2+t^4+4
}
=
t^2+2.
\]

Así:

\[
L
=
\int_0^3
(t^2+2)\,dt.
\]

Por tanto:

\[
\boxed{
L=15
}
\]

\section{Ejercicio 17}
\subsection{Solución}

Sea:

\[
\alpha(t)
=
(t-\sen t,t,1-\cos t),
\qquad
t\in[0,6\pi].
\]

\textbf{a)}

La derivada es:

\[
\alpha'(t)
=
(1-\cos t,1,\sen t).
\]

Como la segunda componente es siempre:

\[
1,
\]

se tiene:

\[
\alpha'(t)\neq(0,0,0)
\]

para todo \(t\).

Por tanto:

\[
\boxed{
\alpha\text{ es regular}
}
\]

\textbf{b)}

La rapidez es:

\[
\|\alpha'(t)\|
=
\sqrt{
(1-\cos t)^2+1+\sen^2t
}.
\]

Simplificando:

\[
\|\alpha'(t)\|
=
\sqrt{
3-2\cos t
}.
\]

Por tanto, la longitud viene dada por:

\[
\boxed{
L
=
\int_0^{6\pi}
\sqrt{
3-2\cos t
}\,dt
}
\]

\section{Ejercicio 18}
\subsection{Solución}

La parametrización es:

\[
\alpha(t)
=
(t^2+1,\,
2+\ln(t+1),\,
e^t).
\]

El punto:

\[
(1,2,1)
\]

corresponde a:

\[
t=0.
\]

La derivada es:

\[
\alpha'(t)
=
\left(
2t,\,
\frac1{t+1},\,
e^t
\right).
\]

Así:

\[
\alpha'(0)
=
(0,1,1).
\]

Por tanto, la recta tangente es:

\[
\boxed{
(x,y,z)
=
(1,2,1)
+
\lambda(0,1,1)
}
\]

\section{Ejercicio 19}
\subsection{Solución}

Sea:

\[
\alpha(t)
=
(t^3,1-t,4t^2).
\]

El vector velocidad es:

\[
v(t)
=
\alpha'(t)
=
(3t^2,-1,8t).
\]

Por tanto:

\[
\boxed{
v(1)=(3,-1,8)
}
\]

El vector aceleración es:

\[
a(t)
=
\alpha''(t)
=
(6t,0,8).
\]

Así:

\[
\boxed{
a(1)=(6,0,8)
}
\]

La celeridad es:

\[
\|v(1)\|
=
\sqrt{
9+1+64
}.
\]

Por tanto:

\[
\boxed{
\|v(1)\|=\sqrt{74}
}
\]

\section{Ejercicio 20}
\subsection{Solución}

Sea:

\[
\alpha(t)
=
(t^2,1+t,1+t^3).
\]

Para:

\[
t=-1,
\]

se obtiene:

\[
\alpha(-1)
=
(1,0,0).
\]

Por tanto:

\[
\boxed{
(1,0,0)\in C
}
\]

El vector velocidad es:

\[
v(t)
=
(2t,1,3t^2).
\]

Así:

\[
\boxed{
v(-1)=(-2,1,3)
}
\]

El vector aceleración es:

\[
a(t)
=
(2,0,6t).
\]

Por tanto:

\[
\boxed{
a(-1)=(2,0,-6)
}
\]

\section{Ejercicio 21}
\subsection{Solución}

La aceleración es:

\[
a(t)
=
(e^t,0,t+1).
\]

Integrando:

\[
v(t)
=
\left(
e^t+C_1,\,
C_2,\,
\frac{t^2}{2}+t+C_3
\right).
\]

Como:

\[
v(0)
=
(1,-3,\sqrt2),
\]

se obtiene:

\[
C_1=0,
\qquad
C_2=-3,
\qquad
C_3=\sqrt2.
\]

Por tanto:

\[
\boxed{
v(t)
=
\left(
e^t,\,
-3,\,
\frac{t^2}{2}+t+\sqrt2
\right)
}
\]

\section{Ejercicio 22}
\subsection{Solución}

La curva es:

\[
\alpha(t)
=
(t,2t,-3t).
\]

Su vector director es:

\[
v=(1,2,-3).
\]

El plano es:

\[
x+2y-3z=10,
\]

cuyo vector normal es:

\[
n=(1,2,-3).
\]

Como:

\[
v=n,
\]

la curva incide perpendicularmente en el plano.

Para hallar el punto de intersección:

\[
t+4t+9t=10.
\]

Por tanto:

\[
t=\frac57.
\]

Así:

\[
\boxed{
P=
\left(
\frac57,\,
\frac{10}{7},\,
-\frac{15}{7}
\right)
}
\]

\section{Ejercicio 23}
\subsection{Solución}

Se sabe que:

\[
\alpha''(t)
=
(0,0,2).
\]

Integrando:

\[
\alpha'(t)
=
(C_1,C_2,2t+C_3).
\]

Como:

\[
\alpha'(0)
=
(2,0,0),
\]

se obtiene:

\[
\alpha'(t)
=
(2,0,2t).
\]

Integrando de nuevo:

\[
\alpha(t)
=
(2t+D_1,D_2,t^2+D_3).
\]

Como:

\[
\alpha(0)
=
(2,0,0),
\]

se obtiene:

\[
\boxed{
\alpha(t)
=
(2t+2,0,t^2)
}
\]

En:

\[
t=1,
\]

\[
\alpha'(1)
=
(2,0,2).
\]

Su norma es:

\[
2\sqrt2.
\]

Por tanto, el vector tangente unitario es:

\[
\boxed{
T(1)
=
\frac1{\sqrt2}(1,0,1)
}
\]

\section{Ejercicio 24}
\subsection{Solución}

La curva es:

\[
\alpha(t)
=
(2\cos t,2\sen t,5t).
\]

El punto:

\[
\left(
0,2,\frac{5\pi}{2}
\right)
\]

corresponde a:

\[
t=\frac{\pi}{2}.
\]

La derivada es:

\[
\alpha'(t)
=
(-2\sen t,2\cos t,5).
\]

Así:

\[
\alpha'\left(
\frac{\pi}{2}
\right)
=
(-2,0,5).
\]

La recta tangente es:

\[
\boxed{
(x,y,z)
=
\left(
0,2,\frac{5\pi}{2}
\right)
+
\lambda(-2,0,5)
}
\]

El plano normal tiene como vector normal:

\[
(-2,0,5).
\]

Por tanto:

\[
\boxed{
-2x
+
5\left(
z-\frac{5\pi}{2}
\right)
=
0
}
\]

\section{Ejercicio 25}
\subsection{Solución}

Sea:

\[
\alpha(t)
=
(t,t^2,t).
\]

El vector velocidad es:

\[
v(t)
=
(1,2t,1)
\]

y la aceleración:

\[
a(t)
=
(0,2,0).
\]

En:

\[
t=1,
\]

se tiene:

\[
v=(1,2,1),
\qquad
a=(0,2,0).
\]

La componente tangencial de la aceleración es la proyección de \(a\) sobre \(v\):

\[
a_T
=
\frac{a\cdot v}{\|v\|^2}v.
\]

Como:

\[
a\cdot v=4,
\qquad
\|v\|^2=6,
\]

se obtiene:

\[
\boxed{
a_T
=
\left(
\frac23,\frac43,\frac23
\right)
}
\]

La componente normal es:

\[
a_N
=
a-a_T.
\]

Por tanto:

\[
\boxed{
a_N
=
\left(
-\frac23,\frac23,-\frac23
\right)
}
\]

\section{Ejercicio 26}
\subsection{Solución}

Para:

\[
D_1
=
\R^2\setminus\{(0,0)\},
\]

existen curvas cerradas que rodean al origen y no pueden contraerse a un punto sin salir del conjunto.

Por tanto:

\[
\boxed{
D_1\text{ no es simplemente conexo}
}
\]

Para:

\[
D_2
=
\R^3\setminus\{(0,0,0)\},
\]

una curva cerrada puede deformarse alrededor del punto eliminado y contraerse sin atravesarlo.

Por tanto:

\[
\boxed{
D_2\text{ es simplemente conexo}
}
\]

\section{Ejercicio 27}
\subsection{Solución}

La superficie de una esfera no contiene agujeros y toda curva cerrada puede contraerse a un punto.

Por tanto:

\[
\boxed{
\text{La esfera es simplemente conexa}
}
\]

En un toro existen curvas cerradas que rodean su agujero y no pueden contraerse a un punto permaneciendo sobre la superficie.

Por tanto:

\[
\boxed{
\text{El toro no es simplemente conexo}
}
\]

\section{Ejercicio 28}
\subsection{Solución}

Sea:

\[
D_1
=
\left\{
(x,y,z)\in\R^3:
1<x^2+y^2+z^2<4
\right\}.
\]

Es una corona esférica tridimensional. Aunque presenta una región interior eliminada, cualquier curva cerrada puede contraerse a un punto dentro del conjunto.

Por tanto:

\[
\boxed{
D_1\text{ es simplemente conexo}
}
\]

Sea:

\[
D_2
=
\left\{
(x,y,z)\in\R^3:
x^2+y^2\neq0
\right\}.
\]

Este conjunto es:

\[
\R^3
\]

sin el eje \(Z\).

Una curva cerrada que rodee al eje \(Z\) no puede contraerse a un punto sin atravesarlo.

Por tanto:

\[
\boxed{
D_2\text{ no es simplemente conexo}
}
\]

\section{Ejercicio 29}
\subsection{Solución}

Sea:

\[
r(u,v)
=
(u,v,u^2+\sen v).
\]

Las derivadas parciales son:

\[
r_u
=
(1,0,2u),
\]

\[
r_v
=
(0,1,\cos v).
\]

Su producto vectorial es:

\[
r_u\times r_v
=
(-2u,-\cos v,1).
\]

La tercera componente es siempre igual a:

\[
1,
\]

por lo que:

\[
r_u\times r_v\neq(0,0,0)
\]

para todo:

\[
(u,v)\in[-2,2]\times[0,2\pi].
\]

Por tanto:

\[
\boxed{
r\text{ es una parametrización regular}
}
\]

\section{Ejercicio 30}
\subsection{Solución}

Sea:

\[
r(u,v)
=
(u+v^2,u+4v,u^2-u).
\]

Las derivadas parciales son:

\[
r_u
=
(1,1,2u-1),
\]

\[
r_v
=
(2v,4,0).
\]

Su producto vectorial es:

\[
r_u\times r_v
=
(4-8u,\,
2v(2u-1),\,
4-2v).
\]

Para:

\[
u=\frac12,
\qquad
v=2,
\]

se obtiene:

\[
r_u\times r_v
=
(0,0,0).
\]

Por tanto:

\[
\boxed{
r\text{ no es una parametrización regular}
}
\]

\end{document}
